%!TEX root = ../../main.tex
\section{결론}
\label{sec:disc-conclusion}

본 논문은 교전이 시작되기 전의 공개 상태에 그 교전의 결과가 얼마나 담겨 있는지를, 경기와 연결된 목표, 곧 동결 평가기 아래 교전 구간의 추정 승률 변화 방향으로 물었다. 한타와 소규모 교전 각각의 안에서, 교전 전 상태의 정규화 로지스틱 모델은 승률과 시간의 스플라인 기준선보다 경기 가중 Brier 점수를 홀드아웃 패치에서 약 \dBrierT{}와 \dBrierS{}만큼 낮추었고, 경기 군집 부트스트랩 구간은 0을 제외하였다. 이득은 대등한 결과의 환원 불가능한 불확실성에 비해 작으며, 그것이 바로 발견이다. 교전 전 공개 상태는 누구나 이미 아는 두 가지, 현재 승률과 경기 시각을 빼고 나면, 교전이 경기를 어디로 움직일지에 관해 측정 가능하지만 소박한 양의 정보를 담고 있다.

검증 스택이 이 주장의 범위를 한정한다. 평가기의 품질은 따로 보고되고, 교전 구간은 짝지은 비교전 구간보다 추정치를 더 크게 움직이며, 라벨은 결정된 물질적·오브젝트 결과와 대응하고, 이득은 작은 변화의 필터에 견딘다. 이득은 한타의 경우 이후 패치와 다른 지역으로의 점수 전용 전이에서 살아남지 못하며, 그때 판별 성분은 여전히 모델에 유리하지만 오보정이 지배한다. 코호트별 모델이 올바른 단위이다. 한타와 소규모 교전을 합쳐도 어느 쪽에도 도움이 되지 않는다. 결과는 모델이 정의한 $\SVI$에 관한 것이며 독립적인 한타 승리 판정이나 인과적 행동 추천을 주장하지 않는다.
